\scanpage{157}
\deleted{où la flèche de droite n'a pas encore reçu droit de cité, dans le formalisme développé jusqu'à présent, du moins}

\deleted{On voit en fait qu'il faut considérer,}
En fait, dans le cas général, si
\[
(255)\qquad\Sigma\subset\Sigma'\deleted{\unsure{\subset X}}
\]
sont des inclusions de parties admissibles de $X$, on veut pouvoir exprimer
\[
\Pi_1\Sigma,\quad\Pi_1(\Sigma'\setminus\Sigma)\quad\deleted{\unsure{et}},\quad
\Pi_1V^*_{\Sigma,\unsure{\Sigma'|X'}},
\]
et \deleted{$\unsure{\Pi_1}$ (avec généralisations convenables)} \inserted{plus généralemt}
\[
\deleted{\unsure{et}}\;\Pi_1V^*_{\Sigma'\setminus\unsure{\Sigma},\unsure{X'|\Sigma'}}
\]
en termes des $\widetilde\Pi_{\Delta_r}$ et des $\Pi_{\Delta_r}$ (et \unsure{sinon|ainsi} introduire les éléments de structure qui manquent).

Plus généralement, considérons
\[
(256)\qquad\deleted{\unsure{\Sigma\subset}}\quad\Sigma'\subset\Sigma''\quad\text{et un }\Sigma_1\subset\Sigma''
\]
\deleted{\unsure{les}} (parties admissibles de $X$) et considérons $V_{\Sigma',\Sigma''}$, considérons la trace de $\Sigma_1$ dessus, qui peut s'écrire
\deleted{$\unsure{\Sigma'\cap}$}
\[
V_{\Sigma',\Sigma''}\cap\Sigma_1=V_{\Sigma'\cap\Sigma_1,\Sigma_1})
\]
\margin{\uline{NB} on ne suppose \unsure{pas} que $\unsure{\Sigma'|\Sigma}\subset\Sigma_1$}
\margin{156}
